\documentclass[fleqn]{tufte-handout}
\title{Speaking about Models}
\author{Hans Halvorson}
\date{June 30, 2023}
% \usepackage{fullpage}
\usepackage{outlines}
\usepackage{amsmath,amssymb}
\usepackage{mathrsfs}
\begin{document}

\maketitle

\begin{enumerate}

\item Most scientists and philosophers take set-theoretic foundations
  of math for granted --- at roughly the level of Halmos' \emph{Naive
    Set Theory}. They assume that the ``niceties'' of set theory won't
  make any difference to their research. And they assume that set
  theory is as safe as anything else we could put in the foundations.

\item Sometimes when we\footnote{especially mathematicians,
    mathematical physicists, and technically oriented philosophers}
  get engaged in detailed investigations, set theory becomes relevant
  --- and we have to decide how far we want to commit to it, and how
  we understand its consequences.

  \begin{itemize}
  \item When applicable results depend on choice axioms
    (e.g. Tarski-Banach, existence of singular distributions)
  \item Putnam's paradox
  \item The hole argument
  \end{itemize}
  
% \item There are alternative foundations of mathematics: constructive
%   foundations, category theory, type theory, HoTT. But if something is
%   not broken, then why try to fix it (especially since we know that no
%   adequate system can prove its own consistency)
  
\item This talk descends from two papers:
  \begin{itemize}
  \item Dimitris Tsementzis and HH, ``Foundations and Philosophy''
  \item HH and JB Manchak, ``Closing the Hole Argument''
  \end{itemize} The first paper considers how replacing set theory
  (ZF) with homotopy type theory (HoTT) might change our philosophical
  logic. What's relevant for us is that HoTT tries to fix some
  problems that might result from a too literal reading of ZF. The
  second paper argues that the hole argument of GTR is based on a
  misunderstanding about how to reason about models (structures in the
  set-theoretic universe) that, once again, results from a too literal
  reading of ZF.

\item The basic puzzle: for models $M$ and $N$ of a theory, how are we
  to understand the situation where $M$ and $N$ are isomorphic but not
  identical?

\item The puzzle restated

  Suppose that $M\neq N$ but that $h:M\to N$ is an isomorphism.

  \begin{itemize}
  \item Defn: $M$ and $N$ are \emph{elementarily equivalent} just in
    case $M\vDash \phi$ iff $N\vDash \phi$.
  \item Fact: If $M$ and $N$ are isomorphic, then $M$ and $N$ are
    elementarily equivalent.
  \item There is a set-theoretic predicate $\Phi$ such that $\Phi (M)$
    but $\neg \Phi (N)$.
  \end{itemize}

\item Example
  \begin{itemize}
  \item Let $T$ be the empty theory in the empty language.
  \item Let $M$ be the model of $T$ with domain $\{ \emptyset \}$.
  \item Let $N$ be the model of $T$ with domain
    $\{ \{ \emptyset \} \}$.
  \item For a $\Sigma$-structure $X$, let $\Phi$ be the set-theoretic
    predicate $\Phi (X)\equiv (\emptyset \in |X|)$. \end{itemize}

\item Is the availability of non-isomorphism invariant predicates a
  bug of set-theoretic foundations, or a feature?
    \begin{itemize}
    \item This fact seems to drive the hole argument in GTR, as well
      as Putnam's paradox.
    \end{itemize}
  \item If it's a bug, can it (and should it) be fixed by changing
    foundations?
    \begin{itemize}
    \item Hermeneutic structuralism versus revolutionary structuralism
    \end{itemize}

\item Putnam's trick: Let $M$ be a model of $T$, and let $h:N\to M$ be
  a bijection of sets. Define $r^N=h^{-1}(r^M)$ for each relation
  symbol $r$. Then the models $M$ and $N$ are isomorphic. But if $h$
  is non-trivial, then $M$ and $N$ disagree on the properties of
  individuals.

  \begin{itemize}
  \item Hilbert, Burgess, Linnebo (hermeneutic structuralism): Pure
    mathematicians are trained to ignore the differences between $M$
    and $N$. It's not the fault of ZF if an untrained person takes $M$
    and $N$ to represent different situations.
  \end{itemize}

\item Does Putnam's trick also lie behind the ``hole argument'' of
  Einstein's general theory of relativity?

  \begin{itemize}
  \item The hole argument is supposed to derive a bad consequence
    (pernicious indeterminism) from a reasonable ontological
    assumption (spacetime is a substance).
  \item Einstein ``Lochbetrachtung'' 1913; Earman and Norton ``What
    price spacetime substantivalism?'' 1987
  \item Simple analogy:
    \begin{itemize}
    \item A model $M$ of $T$ consists of two sets $M_0,M_1$, each of
      which contains two elements, and a functional relation
      $r\subseteq M_0\times M_1$ whose intended meaning is temporal
      development. i.e.\ we stipulate that for each $a\in M_0$ there
      is a unique $a'\in M_1$ such that $\langle a,a'\rangle \in r$.
    \item Given a bijection $h:M_1\to M_1$, the Putnam trick can be
      used to crete another model that is isomorphic to
      $M$.
    \item The theory $T$ thus fails the Montague-Lewis-Earman
      definition of determinism. \end{itemize}
  \end{itemize}
  
\item Benacerraf's antinomy: If $0,1,2,3,\dots $ are sets, then there
  are many facts about them that we have no means of discovering

  \begin{itemize}
  \item Is $\emptyset\in 1$?
  \end{itemize}
  
\item A middle way?

  \begin{itemize}
  \item There are too many set-theoretic predicates: they detect
    differences that are \emph{not} really there
  \item There are too few $\Sigma$-sentences: they cannot detect real
    differences
    \begin{itemize}
    
    \item Example: $(\mathbb{R},\leq )$ and $(\mathbb{Q},\leq )$ are
      elementarily equivalent.
    \item Recall: In the case of finite structures, elementary
      equivalence implies isomorphism.
    \end{itemize}
  \end{itemize}

\item Open question: Can we explicate mathematicians' talk about
  ``properties invariant under isomorphism'' in terms of definable
  properties of models? This would suggest a \emph{preferred language
    for speaking about models of a theory}.
  \begin{itemize}
  \item Topological properties: Hausdorff, compact, metrizable, simply
    connected, etc.
  \item Group properties: abelian, number of subgroups, number of
    normal subgroups, etc.
  \item Relativistic spacetimes: singular, globally hyperbolic,
    etc.
    \begin{itemize}
    \item The issue is of real importance here. E.g. it is easy to
      mistake coordinate singularities for genuine singularities. In
      fact, it's not so easy to check that proposed definitions of
      ``is singular'' are invariant under
      isomorphism \end{itemize} \end{itemize}


\item Some first ideas
  \begin{itemize}
  \item Properties of models that are representable by first-order
    sentences
    \begin{itemize}
    \item Partial orders: has a left endpoint
    \item Boolean algebras: atomic, atomless
    \item Groups: abelian, torsion free, divisible, nilpotent,
      solvable
    \end{itemize}
  \item There are invariant properties of models that are not
    representable by first-order sentences
    \begin{itemize}
    \item Partial orders: any two elements are connected by a finite
      chain
    \item Peano arithmetic: is an end extension
    \item Groups: has $n$ distinct normal subgroups, finite, free,
      simple
    \end{itemize}
  \end{itemize}

\item Where can we get more insight about how to describe the
  invariant properties of a model $M$?
  \begin{itemize}
  \item The category $D(M)$ of definable subsets and definable
    functions.\footnote{see Pillay, ``Model theory'' \emph{Notices of
        the AMS} 2000}
    \begin{itemize}
      \item Fact: If $h:M\to N$ is an isomorphism, then $D(M)$ is
        isomorphic to $D(N)$.
      \item Conjecture: If $D(M)$ is isomorphic to $D(N)$, then $M$ is
        isomorphic to $N$.
    \end{itemize}
  \item McGee has results that indicate that all the invariant
    features of $M$ can be described in $\mathscr{L}_{\infty\infty}$.
  \item Vaananen has a proposal about how the signature $\Sigma$
    defines higher order sorts.\footnote{``Sort logic and foundations
      of mathematics.'' In \emph{Infinity and truth}}
  \end{itemize}

\item Revolutionary Structuralism I: Elementary theory of the category
  of sets (ETCS)

  \begin{itemize}
  \item Key fact: In ETCS, no two sets can share a point in common
  \item ETCS is a two sorted first-order theory: objects and arrows
  \item Axioms: initial object, terminal object, Cartesian products,
    coproducts, natural number object, etc.
  \item A \textbf{point} of $X$ is a morphism $p:1\to X$
  \item Each isomorphism $h:X\to Y$ gives a way of identifying points
    in $X$ with points in $Y$.
  \end{itemize}

\item Mitchell's Theorem: ETCS is ``equivalent'' to a weak form of ZF
  set theory

  \begin{itemize}
  \item Open question: what does ``equivalence'' mean here?
  \begin{itemize}
  \item If it's merely ``equal logical strength'' then it's too weak
    to carry interpretive weight (e.g.\ G\"odel translation of
    intuitionistic logic)
  \item Bi-interpretability? Definitional equivalence?
  \item Morita equivalence?
  \end{itemize}
\item Could category-theoretic foundations could be \emph{equivalent}
  in some precise sense to set-theoretic foundations?
\end{itemize}

\end{enumerate}



\end{document}

%%% Local Variables:
%%% mode: latex
%%% TeX-master: t
%%% End:

\item Example: Let $T$ be the empty theory in the empty
  language. There is a family $\{ M_\kappa \}$ of models of $T$ and a
  family $\Phi _k$ of set-theoretic predicates such that:
  \begin{enumerate}
  \item $|M_\kappa |=1$ for all $\kappa$;
  \item $M_\kappa$ is isomorphic to $M_\lambda$ for all
    $\kappa ,\lambda$;
  \item $\Phi _\kappa (M_\lambda )$ iff $\kappa =\lambda$.
  \end{enumerate}

